%=====================================================
%
% Variables (same scheme as the sample document)
%
%=====================================================

\def\doctype{article}

\def\emptypages{false}
\def\showparts{false}

\def\ttitlepage{titlepage.tex}
\def\ttitlepagetitletop{Local Coding Agents on One Workstation, Second Night}
\def\ttitlepagesubtitletop{What the Serving Stack, the Prompt and the Model Each Contribute}
\def\ttitlepagetitlemiddle{Lab Report}
\def\ttitlepagesubtitlemiddle{Paired measurements of 15--16 October 2026}

\def\torganization{Independent Research}
\def\tsubmissiondate{16 October 2026}
\def\tversion{0.1 (draft)}
\def\tauthor{Matthias Nott}
\def\tcontact{~}
\def\tcontactid{~}

\def\tfooterleftlabelup{Title}
\def\tfooterleftlabeldown{Author}
\def\tfooterleftup{Local Coding Agents, Night 2}
\def\tfooterleftdown{Matthias Nott}
\def\tfootermiddleup{Serving, Prompt, Model}
\def\tfootermiddledown{Lab Report}
\def\tfooterrightlabelup{Date}
\def\tfooterrightlabeldown{Version}
\def\tfooterrightup{October 2026}
\def\tfooterrightdown{0.1}

\def\showtitlepage{true}
\def\showlistofs{true}
\def\showtoc{true}
\def\showlot{true}
\def\showlof{false}
\def\showloe{false}
\def\showlod{false}

\def\showreferences{true}
\def\showappendix{false}
\def\showendnotes{false}
\def\showanyglossary{false}
\def\showacronyms{false}
\def\showglossary{false}
\def\showstatic{false}
\def\showindex{false}

\def\showpageheaders{true}
\def\showpagefooters{true}
\def\showwordcount{false}
\def\showwordcountback{false}

\input{templates/cfg/zz_layout}

\setcounter{tocdepth}{2}

\usepackage[
  style=apa
, citestyle=authoryear
, backend=biber
, maxnames=3
, minnames=1
, sorting=nyt
, sortcites=true
, block=space
, safeinputenc
, natbib=true
, backref=true
, uniquename=init]{biblatex}
\DeclareLanguageMapping{american}{american-apa}
\addbibresource{Bibliography.bib}
\input{templates/cfg/zz_config_bib}

\usepackage{indentfirst}

\newcommand{\cmd}[1]{\texttt{#1}}

% Alternate row colours in every table (xcolor is loaded with the table option by the layout)
% Breathing room between tables and the running text
\setlength{\intextsep}{16pt plus 2pt minus 2pt}
\setlength{\textfloatsep}{18pt plus 2pt minus 4pt}
\setlength{\floatsep}{16pt plus 2pt minus 2pt}

\definecolor{rowtint}{gray}{0.93}
\AtBeginEnvironment{table}{\global\csname @rowcolorstrue\endcsname\rowcolors{2}{rowtint}{white}}
\newcolumntype{P}[1]{>{\raggedright\arraybackslash}p{#1}}
\setcounter{secnumdepth}{2}
\newcommand{\todo}[1]{\textcolor{red}{[TODO: #1]}}

\begin{document}

%=====================================================
% Title page and abstract
%=====================================================

\input{tex/\ttitlepage}

\begin{center}
\begin{minipage}{0.86\textwidth}
\section*{Abstract}

The first night showed that open-weight models on one workstation could not yet replace a cloud worker on real
application fixes, but its comparisons were run in blocks and its local servers were still defective. This second
night repeats the question with paired arms: the two arms of a comparison start together on the same task, only one
model server is loaded at a time, and the local runs are graded by the hidden tests of the fixes that actually landed.
Three findings stand out. First, naming the file and the function (the \emph{precise} prompt) lifted the local
Qwen3-Coder-30B-A3B worker from 1 of 9 to 5 of 9 passes against the orchestrator's specification; all four discordant
pairs favoured the precise prompt (exact sign test $p=0.125$, two-sided, which is the smallest value nine pairs can
give with four discordant). Second, serving still decided the outcome: a second model family ran through a multimodal
code path that never hit the prefix cache, so every turn re-read 40--50k tokens; after the fix, a smoke run fell from
55--61\,s to 23--37\,s and a one-function fix ran within about twice the time of the cloud reference. Third, a cloud
plan written from the parent commit let three different local models pass a task that none of eight local precise or specification runs
passed. A hardware watchdog reset at 23:43, with about 117\,GB of model servers resident on a 128\,GB machine, set the
constraint for the rest of the night: one model server at a time. The samples are small and every statement is
conditional on them. The final ranking is Sonnet~5.5 (4 of 4) ahead of GLM-5.3-flash and GLM-5.3 (3 of 4 each), Haiku~4.5 (2 of 4)
and then the local models, of which Ornith-1.5-35B-A3B and Qwen3-Coder-30B-A3B are level on passes (7 of 16 and 4 of 8)
and the dense Qwen3.8-27B is too slow to use; no local model passed the three-file task.

\end{minipage}
\end{center}


\clearpage

%=====================================================
% Table of contents and list of tables
%=====================================================

\sequals{\showpageheaders}{true}{
	\pagestyle{fancy}
	\fancyhead{}
	\fancyhead[RO,LE]{\thepage}
	\fancyhead[LO]{\slshape CONTENTS}
	\fancyhead[RE]{\slshape CONTENTS}
}{}
\sequals{\showpagefooters}{true}{\fancyfoot[C]{}}{}
\pdfbookmark[-1]{\contentsname}{toc}
\tableofcontents
\listoftables
\clearpage

\sequals{\showpageheaders}{true}{
	\pagestyle{fancy}
	\fancyhead{}
	\fancyhead[LE,RO]{\thepage}
	\fancyhead[RE]{\slshape \rightmark}
	\fancyhead[LO]{\slshape \rightmark}
}{}
\sequals{\showpagefooters}{true}{\fancyfoot[C]{\pfooter}}{}
\renewcommand{\sectionmark}[1]{\markboth{\thesection.\ #1}{\thesection.\ #1}}
\renewcommand{\subsectionmark}[1]{}

%=====================================================
% Body
%=====================================================

\section{Introduction}\label{sec:intro}

The first report measured whether open-weight models on a single 128\,GB workstation could act as
the workers in an agentic coding workflow \citep{Anthropic-2026aa}. Its answer was no for real application fixes, but
it left three questions open: whether the failures were a matter of finding the code or of changing it, whether the
local servers were themselves the bottleneck, and whether the blocked order of the runs had biased the comparison.

The question for this report is narrower and more practical: \emph{can a local coding setup run beside normal work?}
That means one model server of moderate size, no extra hardware, and a worker that completes small, well-localised
fixes in a time the operator will accept, while the machine is still used for other things. The cloud models are
yardsticks only; the aim is a local solution, in speed and in quality.

Each experiment therefore answers one of three questions: which \emph{serving setup}, which \emph{prompt}, or which
\emph{model} makes a local worker pass more often and faster. Section~\ref{sec:changes} describes what was changed
since the first night, Section~\ref{sec:setting} the setting, Section~\ref{sec:method} the method, and
Section~\ref{sec:results} the results; Section~\ref{sec:reboot} reports the reboot that constrained the night.

\section{What Changed Since Night 1}\label{sec:changes}

\begin{description}
	\item[Paired arms.] The two arms of a comparison run on the same task and start together, so thermal state, load and
	the build queue affect both alike. Night 1 ran its arms in blocks.
	\item[One model server at a time.] After the reboot of Section~\ref{sec:reboot}, models are compared phase by
	phase on an otherwise idle machine; within a phase, the paired arms share the server.
	\item[Path repair.] Local models repeatedly mistyped absolute paths into the arm's own clone, one character off from
	the working directory, and received ``file does not exist'' again and again. A worker-only hook now repairs a
	non-existent absolute path when its prefix is within three edits of the working directory and the remainder exists
	below it. Arms started after the hook was registered use it; both arms of a pair start together, so pairs remain
	comparable.
	\item[Isolation.] As on night 1, each run starts in its own clone with only the parent commit's history; hidden tests
	are overlaid at grading time. Every row is tagged with the night so the two nights stay apart.
\end{description}

\section{Setting}\label{sec:setting}

All measurements ran on one Apple MacBook Pro with an M5 Max and 128\,GB of unified memory, set to maximum
performance by the operator; the thermal level and the load average were logged with every pair. Memory pressure and
swap use were logged every 10\,s during runs. Local models were served with vllm-mlx \citep{vllm-mlx-2026aa} through
patched launchers. Workers were started through PAI \citep{Nott-2026aa}; sessions exchanged messages through AIBroker
\citep{Nott-2026ab}. The workloads are anonymised: Project A is the native PDF application that supplies the coding
tasks, Project B the document-management system that shares the machine.

\begin{table}[htbp]
\centering
\small
\caption{Models in the second night}\label{tab:models2}
\begin{tabular}{P{4.2cm}P{2.2cm}P{2.6cm}P{5.0cm}}
\toprule
Model & Placement & Size & Role \\
\midrule
Qwen3-Coder-30B-A3B \citep{Qwen-Team-2025aa} & local & 4-bit & Routed local model (``local30'') \\
Ornith-1.5-35B-A3B & local & 4-bit & New local candidate \\
Qwen3.8-27B & local, dense & 6-bit & New local candidate (``qwen38'') \\
GLM-4.7-Flash \citep{Zhipu-AI-2026aa} & local & 6-bit & New local candidate (``glm47'') \\
Claude Sonnet 5.5 \citep{Anthropic-2026ab} & cloud & -- & Yardstick \\
Claude Haiku 4.5 & cloud & -- & Yardstick \\
GLM-5.3 & cloud & -- & Yardstick \\
GLM-5.3-flash & cloud & -- & Yardstick \\
\bottomrule
\end{tabular}
\end{table}

\section{Method}\label{sec:method}

\subsection{Tasks}

The tasks are the replay tasks of night 1: fixes that already landed in Project A, started from the parent commit and
graded by the hidden tests that came with the fix. Four tasks have a precise prompt and were used tonight:
\emph{log-size} (one function, +21 lines in the landed fix), \emph{box-press} (2 files, +79), \emph{view-decimals}
(1 file, +6) and \emph{save-model-ref} (3 files, +64).

\subsection{Three prompts, and a plan}

The \emph{task} prompt is the maintainer's symptom in their own words. The \emph{spec} prompt is what the orchestrator
wrote when it delegated the fix: the diagnosis and the owning part of the code, without the fix. The \emph{precise}
prompt names the exact file and function. The \emph{plan} prompt adds an edit plan (files, functions, what to change,
no code) written by a cloud model from the parent commit without access to the fix: ``cloud head, local hands''. A leak
audit checked that no prompt contained the fix.

\subsection{Grading and pairing}

A patch passes when the result builds and the hidden tests pass; no model judges another. Wall time is censored at the
task deadline. Pairs are two arms of the same task started together; discordant pairs are counted and an exact sign
test is applied. Wilson intervals \citep{Wilson-1927aa} are given for single arms. Tonight's samples are one to three
repeats per task and cell.

\section{Results}\label{sec:results}

\subsection{Precise against specification on one local model}\label{sec:e1}

Table~\ref{tab:e1} gives the nine completed pairs for the local30 model. Precise passed 5 of 9 (Wilson 27--81\,\%),
the specification 1 of 9 (2--44\,\%). Of the nine pairs, four were discordant, all in favour of precise; the exact sign
test gives $p=0.125$ two-sided, which is the smallest value attainable with four discordant pairs. The two tasks
log-size and box-press, where the precise prompt names the function to change, account for all the gain; precise runs also used about half the
turns. view-decimals and save-model-ref failed in both arms in every repeat: for view-decimals the worker reached the
right file and the change itself was wrong, so those failures are about making the change, not finding it.

\begin{table}[htbp]
\centering
\small
\caption{local30, specification against precise prompt (paired; minutes / turns where recorded)}\label{tab:e1}
\begin{tabular}{lcc}
\toprule
Pair & Spec & Precise \\
\midrule
log-size, rep 1 & fail 14.6 / 41 & \textbf{pass} 11.6 / 27 \\
box-press, rep 1 & fail 14.7 / 37 & \textbf{pass} 15.2 / 33 \\
view-decimals, rep 1 & fail 6.2 / 40 & fail 2.2 / 15 \\
save-model-ref, rep 1 & fail 18.4 & fail 12.9 \\
log-size, rep 2 & fail 13.0 / 79 & \textbf{pass} 6.6 / 27 \\
box-press, rep 2 & fail 15.5 / 52 & \textbf{pass} 15.2 / 32 \\
view-decimals, rep 2 & fail 5.8 / 35 & fail 13.4 / 53 \\
save-model-ref, rep 2 & fail 14.2 / 35 & fail 14.1 / 31 \\
log-size, rep 3 & \textbf{pass} 8.8 / 40 & pass 6.2 / 28 \\
\midrule
Passed & 1 / 9 & 5 / 9 \\
\bottomrule
\end{tabular}
\end{table}

\subsection{Cloud yardsticks}

Table~\ref{tab:cloud} gives the cloud models on the precise prompt, one agent at a time, on an otherwise idle machine
(the GLM-5.3 run was made on a loaded machine). save-model-ref, which changes three files and adds 64 lines, is the
task that separates Sonnet 5.5 from all others; Sonnet takes 6 to 7 minutes on the two tasks that no local model has
passed. On the specification prompt GLM-5.3 passed 3 of 4 as well (log-size, box-press and, after 39.6 minutes and 117
turns, save-model-ref, graded by hand), failing view-decimals; with precise it passed view-decimals and lost
save-model-ref.

\begin{table}[htbp]
\centering
\small
\caption{Precise prompt, cloud yardsticks (minutes to finish)}\label{tab:cloud}
\begin{tabular}{lccccc}
\toprule
Model & view-decimals & save-model-ref & log-size & box-press & Passed (median min) \\
\midrule
Sonnet 5.5 & pass 6.0 & pass 7.0 & pass 2.0 & pass 5.2 & 4/4 (5.6) \\
GLM-5.3 & pass 13.7 & fail 21.7 & pass 3.9 & pass 19.8 & 3/4 (16.8) \\
GLM-5.3-flash & pass 13.9 & fail 40.0 & pass 4.9 & pass 8.3 & 3/4 (11.1) \\
Haiku 4.5 & fail 8.3 & fail 10.6 & pass 3.0 & pass 8.1 & 2/4 (8.2) \\
\bottomrule
\end{tabular}
\end{table}

\subsection{Local models, one at a time}

Table~\ref{tab:local} compares the local models on the precise and plan prompts, each model served alone with the
paired arms sharing the server; cells show the result with minutes (turns where recorded). Ornith-1.5-35B-A3B was run
in two rounds: round 1 without the levers of Section~\ref{sec:failures}, round 2 with the loop breaker, the
first-turn retry and the final path hook. Rows of the first, defective serving (no prefix cache) and the two repeat-2 arms
killed at the reorder stop are excluded. Ornith ran with thinking off, which is the configuration that is fast enough to
use (Section~\ref{sec:serving}). Over both rounds Ornith passed 3 of 8 with the precise prompt and 4 of 8 with the plan
prompt; Qwen3-Coder-30B (local30) passed 2 of 4 with each. Plan wins view-decimals for every model that ran it (4 of 4
graded arms pass; Ornith twice, local30, Qwen3.8), precise wins box-press for local30 and Ornith. All eight local
precise or specification arms on view-decimals failed (local30 four, Ornith two, local30 phase one, Qwen3.8 one), so
the localisation hypothesis does not extend to that task. On save-model-ref no local arm passed (0 of 12).

\begin{table}[htbp]
\centering
\small
\caption{Local models by prompt (night 2; minutes, turns where recorded; 20 or 40 min is the deadline; passed with median wall minutes of graded arms)}\label{tab:local}
\resizebox{\textwidth}{!}{\begin{tabular}{llccccc}
\toprule
Model & Prompt & view-decimals & save-model-ref & log-size & box-press & Passed (median) \\
\midrule
Ornith round 1 & precise & fail 20.0 & fail 40.0 & pass 4.6 / 27 & fail 0.1 (empty first turn) & 1/4 (12.3) \\
Ornith round 1 & plan & \textbf{pass} 20.0 & fail 40.0 (no-op loop) & pass 9.5 / 49 & fail 25.0 & 2/4 (22.5) \\
Ornith with levers & precise & fail 20.0 & fail 40.0 & pass 4.9 / 25 & pass 14.0 / 94 & 2/4 (17.0) \\
Ornith with levers & plan & \textbf{pass} 20.0 & fail 40.0 & pass 9.9 / 55 & fail 12.5 / 71 & 2/4 (16.3) \\
local30 & precise & fail 2.6 / 27 & fail 31.2 / 52 & pass 6.0 / 26 & pass 13.3 / 36 & 2/4 (9.7) \\
local30 & plan & \textbf{pass} 10.7 / 36 & fail 40.0 & pass 9.6 / 29 & fail 18.2 / 60 & 2/4 (14.5) \\
Qwen3.8-27B & precise & fail 20.0 & fail 40.0 & pass 15.1 & void: pane gate (06:51) & 1/3 \\
Qwen3.8-27B & plan & \textbf{pass} 20.0 & fail 40.0 & pass 15.0 & void: pane gate (06:51) & 2/3 \\
GLM-4.7-Flash & precise & void: pane gate & void: pane gate & void: pane gate\textsuperscript{a} & void: pane gate & -- \\
\bottomrule
\end{tabular}}

\smallskip\footnotesize
Void: the arm was refused by the worker pane gate (the selected terminal tab belonged to another session) and produced no model result. For GLM-4.7-Flash the server was up at 06:53, but all 8 arms were refused. \textsuperscript{a}A log-size pass of 15.0 minutes exists for GLM-4.7-Flash from 23:18, before the 23:43 reboot, with four servers loaded; pre-reboot, four servers loaded, not comparable, and not counted.
\end{table}

On log-size, the one-function fix, Ornith with the repaired server passed the precise prompt in 4.6 minutes and 27
turns, against 2.0 minutes and 12 turns for Sonnet and 6.2 to 11.6 minutes for local30 on night 1: a local model is
within about a factor of two of the cloud reference on a one-function fix.

\subsection{Serving findings}\label{sec:serving}

Serving decided the outcome again. Table~\ref{tab:serving2} lists the measured defects and what each fix changed.

\begin{description}
	\item[The multimodal path.] Ornith's configuration contains a vision section, so vllm-mlx loaded it through the
	multimodal batch generator, which never hits the prefix cache; the existing text-forcing patch missed the batched
	engine. Every turn logged 40--50k remaining prompt tokens with nothing cached, so turns took 15--100\,s for 130--2000
	output tokens (effective 8--22 tokens/s against 52.8 tokens/s for a single request). Both view-decimals arms ran to
	the 20-minute deadline without an edit; those rows measure the server, not the model, and are marked as such.
	\item[The fold launcher.] Two of the new templates reject a system message after the start, so the launcher folds
	system messages. It appended each turn's trailing system reminder to the last message; the client does not resend
	it, so the history changed every turn and the cache missed. Keeping only the leading system message fixed it.
	\item[Thinking.] With thinking on, the chat template drops earlier reasoning from the history, so every turn misses
	the cache (tested). Thinking was therefore turned off for Ornith and Qwen3.8.
	\item[Output cap.] Requests carried a maximum of 64{,}000 output tokens; the new providers now cap output at 8192, as
	local30 does.
\end{description}

\begin{table}[htbp]
\centering
\small
\caption{Serving defects of night 2, before and after}\label{tab:serving2}
\begin{tabular}{P{4.0cm}P{4.5cm}P{4.5cm}}
\toprule
Defect & Before & After \\
\midrule
Multimodal path, fold launcher, thinking (Ornith, Qwen3.8) & Prefix cache never hit; 15--100\,s per turn; smoke run 55--61\,s, 2.1--4.5\,s per turn & Cache 11 hit / 1 miss, about 95\,\% of each prompt cached; smoke run 23--37\,s, 0.8--1.5\,s per turn \\
Single-turn speed, shared machine (not a defect) & Ornith: prefill 1679, decode 52.8 tokens/s & Qwen3.8-27B dense: prefill 560, decode 15.2 tokens/s \\
\bottomrule
\end{tabular}
\end{table}

The three causes were fixed together before the next run, so the table gives their combined effect; separate before and
after numbers for each cause were not measured. A fourth fault appeared late and was not fixed: after an aborted
request, the Qwen3.8 server rescheduled a 54{,}445-token prompt in full with no cached prefix, which cost minutes per
turn for the dense model (one serving fault more to fix before a dense model is usable).

\subsection{Failure classes and levers}\label{sec:failures}

Three failure classes recur in the local runs.
\begin{description}
	\item[Path typo.] Absolute paths into the arm's own clone mistyped by a few characters. Lever: the path-repair hook
	(Section~\ref{sec:changes}). The repair log counts 20 repairs in four Ornith arms (log-size and box-press, round 2), 12 in five Qwen3.8 arms, and none
in a cloud or local30 arm; these counts are for arms started after the final hook was registered, and the sample is
too small to compare rates between providers.
	\item[Repetition loop.] Ornith on save-model-ref sent the same Edit with identical old and new strings six or more
	times in a row (``no changes to make''), every 12\,s; the plan arm of the same task looped on a no-op edit until the
	deadline. Candidate lever: a hook that, after $N$ identical failing calls, returns a stronger message to re-read the
	file and change approach. It was not applied during the A/B runs.
	\item[Empty first turn.] Ornith on box-press with the precise prompt answered with five tokens on turn 1 and ended
	without a tool call; the run was over in six seconds and counts as a fail. Candidate lever: retry in the runner on an
	empty first turn.
\end{description}
With the levers (Table~\ref{tab:local}), Ornith passed the same number as before (precise 2 of 4 against 1 of 4, plan 2
of 4 against 2 of 4). On save-model-ref the failure changed from a no-op loop without a build to a build with wrong
behaviour, and no pass was gained. The one new pass, box-press with the precise prompt (14.0 minutes, 94 turns), was one
where no lever fired (no nudge or denial in its transcript), so it is run-to-run variance and not an effect of the levers.

\section{The Reboot}\label{sec:reboot}

At 23:43 the machine was reset by the hardware watchdog (boot fault record ``wdog, reset\_in\_1''; no kernel panic).
The window server's graphics submit thread had been blocked in the kernel graphics driver for 40\,s and then 80\,s in
every one of 12 samples. Four MLX servers were resident, of 33.5, 31.5, 29.3 and 22.6\,GB, about 117\,GB on a 128\,GB
machine, with a further 8.6\,GB for a document application and 8.2\,GB for a virtual machine; the compressor and swap
threads were active, but no process was killed, and the thermal level had returned to 0 ten minutes earlier. Our
reading is that GPU contention from four servers and about six agents, plus builds, starved the display server, with
memory pressure aggravating it. We did not isolate the cause further.

The constraint that followed is itself a finding: on one Apple GPU, the number of servers times the number of
concurrent agents is the hard limit, not the amount of RAM. For the rest of the night: at most two servers loaded and
later one, weights plus cache under about 75\,GB and finally about 35\,GB, and at most three local agents at once.

\section{Discussion}\label{sec:discussion}

\textbf{The prompt matters most for localisable fixes.} Naming the file and the function turned two of the four
tasks from consistent failure into passes, at about half the turns, which supports the localisation
hypothesis of night 1 for those tasks. It did not help where the change itself was hard (view-decimals) or spread over
files (save-model-ref). The cloud plan helped on view-decimals for two models, so a divided workflow, cloud head and
local hands, is the more promising direction; with these sample sizes it is a hypothesis, not a result.

\textbf{Serving is a precondition.} A model that looks unable to finish may be waiting on re-reading its prompt. Two
Ornith arms were lost to a cache defect that was invisible in the pass-fail table. Every future model should be
smoke-tested for cache hits before being judged.

\textbf{Failures are partly mechanical.} The path typo, the loop and the empty first turn are deterministic and
cheap to guard against, unlike a wrong understanding of the code. What the levers recovered was a failure mode, not passes: no cell changed from fail to pass because of a lever.

\section{Recommendation}\label{sec:recommendation}

For local coding beside normal work we recommend Ornith-1.5-35B-A3B at 4 bits, served through the text-path fold launcher
with thinking off and the prefix cache on, with one model server and at most two agents at a time, and the
path-confinement hook on. Evidence: it was the only local model fast enough for the loop (0.8 to 1.5 s per turn after the
cache fix, against up to 190 s for a single Qwen3.8 turn), and it passed one-function fixes within about twice Sonnet's
time (log-size 4.6 to 4.9 minutes against 2.0). Weakness: it is level with Qwen3-Coder-30B on passes (7 of 16 against
4 of 8), with one to four runs per cell.

Give it a cloud-written edit plan (cloud head, local hands) for one- and two-function fixes. Evidence: the plan prompt
passed view-decimals in 4 of 4 graded arms on three local models, where 0 of 8 precise or specification arms passed.
Weakness: one task, and the plan made box-press worse for two models (changes of 43 and 80 lines against 5), so the plan
helps only where the change is hard to derive and hurts where it invites a broad edit.

Keep multi-file work, the save-model-ref class, on Sonnet~5.5: it passed in 7.0 minutes, while no local arm passed in
12 tries and GLM-5.3 and GLM-5.3-flash failed it under the precise prompt (n=1 each). Do not use the dense Qwen3.8-27B:
0.9 to 8.1 effective tokens per second with two agents, plus the abort fault above. Haiku~4.5 (2 of 4) fell below
GLM-5.3-flash (3 of 4) on these tasks, so GLM-5.3-flash is the cheap cloud fallback (median 11.1 minutes against 8.2 for
Haiku); the gap is one task and not significant.

\section{Limitations}\label{sec:limits}

The samples are small: nine pairs for the central comparison, one to three repeats per cell elsewhere, and a single
repeat for Ornith. The sign test cannot reach conventional significance with four discordant pairs. The four tasks come
from one subsystem of one code base, and three of them are one-function fixes, which favours the precise prompt. Ornith
and Qwen3.8 ran with thinking off, so their results do not describe the models at their best. The GLM-5.3 yardstick
ran on a loaded machine; wall times of cloud arms depend on the provider. Runs were censored at the deadline. Prompts
and plans were written by Claude models, and a Claude model is also a contestant; grading is mechanical, which limits
but does not remove this. Finally, the arms of a pair share one server and the machine's memory, so the absolute
speeds are those of a shared server and not of a dedicated one.

\clearpage
\markboth{\MakeUppercase{\slshape References}}{}
\phantomsection
\addcontentsline{toc}{section}{References}
\printbibliography[title={References}]

\end{document}
